\documentclass[a4paper,10pt]{article}
\newcommand\subdir{/home/koen/LaTeX-setup/}
\input{\subdir/LaTeX-files/imports.tex}
\input{\subdir/LaTeX-files/master-internship.tex}
\usepackage{\subdir/LaTeX-files/aastex}
\usepackage{natbib}
\bibliographystyle{\subdir/LaTeX-files/astroadstitle.bst}

%Set the title, Author and date
\title{Notes Week 13}
\author{Koen Essers}
\tutorial{Master Internship}
\date{\today}

%Set the header style
\header{\tutorialstring}

%Begin the document
\begin{document}
\setuplayout
\section{Poster}
I spent the majority of this week on designing the poster for the NAC.

\section{Mistake}
Even though I looked at the corrigendum from \citet{temmink2025}, and I explicitly set \lstinline[style=output, breaklines=true]|use_radiation_corrected_transfer_rate| in my models, I very stupidly set it to \lstinline[style=output, breaklines=true]|.true.| instead of \lstinline[style=output, breaklines=true]|.false.|. This is fixed from now on.


\section{Carbon transfer}

I am interested in seeing how much carbon gets transferred to the accretor.
\begin{figure}[ht]
    \marginfignote{0}{This plot shows the mass of \(\text{C}_{12}\) that is transferred at the end of the evolution of the binary. I compare it with the total amount of \(\text{C}_\text{12}\) available in the star as a function of time. The binary models all transfer a significant part of the available \(C_\text{12}\) to the companion, but not the total amount. This must be due to wind mass loss, and also for these models, the radiation corrected transfer.}
    \centering
    \incplot{w13-transfer-1}
\end{figure}

\begin{figure}[ht]
    \marginfignote{0}{This plot shows the amount of \(\text{C}_{12}\) that is left in the star for the various models. It clearly drops significantly when mass transfer occurs, which is what we would expect.}
    \centering
    \incplot{w13-transfer-2}
\end{figure}

By assuming the carbon enriched mass that gets accretor onto the accretor is effectively mixed before the observations, we can quite easily compute the carbon abundance of the accretor after mass transfer.

We use that
\begin{align*}
    m (\text{C}_{12})_\text{final} &= m (\text{C}_{12})_\text{initial} + m (\text{C}_{12})_\text{transferred}\\
    X(\text{C}_{12}) &= \frac{m (\text{C}_{12})_\text{final}}{M_\text{star}}\\
\end{align*}

I also assume that the mass fraction of \(\text{C}_{12}\) in the accretor star is constant and does not change with the mass of the initial accretor\sidenote{Such that \(X(\text{C}_{12})_\text{heavy star} = X(\text{C}_{12})_\text{light star}\)}. 

The plot below shows the approximate carbon 12 abundance of the accretor during the evolution.
\begin{figure}[ht]
    \centering
    \incplot{w13-transfer-3}
\end{figure}

For models that interact early on during the TPAGB, the abundance drops, but once the carbon abundance in the TPAGB star exceeds the abundance of the accretor, the carbon abundance of the accretor will grow after mass transfer. The thick lines show the carbon abundance of the donor.

To compute the C/O-number ratio, we need to compute the mass ratio of \(\text{O}_\text{16}\) as well. Unfortunately, this element changes in abundance during the TPAGB:

\begin{figure}[ht]
    \marginfignote{0}{The \(\text{O}_{16}\) mass fraction evolves quite similar to that of \(\text{C}_{12}\), just a little less extreme, which causes the CO-number ratio to shift for single TPAGB stars.}
    \centering
    \incplot{w13-o16}
\end{figure}

This means we cannot assume it to remain constant and we must use the same procedure to compute the mass fraction of \(\text{O}_{16}\) for the accretor.

If we do this, and divide the two and multiply by the numerical factor \(16 / 12\) to turn the faction into a number ratio, we obtain:

\begin{figure}[ht]
    \centering
    \incplot{w13-transfer-4}
\end{figure}

Under these assumptions it seems like only if a star interacts quite late in the evolution, and with a significantly low mass ratio, it can appear as a carbon enhanced star. The assumptions are on one hand quite conservative. We assume the star mixes the accreted material throughout the \emph{whole} star, which significantly reduces the effect of the accretion on the surface abundance compared to assuming no mixing. On the other hand, we assume \sidenote{nearly, because of the radiation losses} all of the mass is accreted, which is very probably an overestimation of the amount of total transferred mass.

As a final plot, I show the final C/O-number ratio for the accretor for the whole grid:

\begin{figure}[ht]
    \marginfignote{0}{This plot nicely shows where the models lie that achieve a C/O-number ratio greater than 1, and where the model lie that achieve ratios lower than 1. It aligns well with what a wrote above.}
    \centering
    \incplot{w13-co-1}
\end{figure}



\newsection{Circumbinary toroid}
I looked through MESA's source code for its implementation of the circumbinary coplanar toroid. It is based on the same 2 parameters as \citet{soberman1997}, where \(\delta \) determines how much of the \emph{transferred} mass is lost to a toroid. Specifically, they define \lstinline[style=output, breaklines=true]|mdot_system_cct = b% mtransfer_rate * b% mass_transfer_delta|. This is set as follows:
\begin{minted}[fontsize=\small, breaklines]{fortran}
! $MESA_DIR/binary/private/binary_private_def.f90
! subroutine adjust_mdots

! mdot_system_transfer is mass lost from the vicinity of each star
! due to inefficient rlof mass transfer, mdot_system_cct is mass lost
! from a circumbinary coplanar toroid.
if (b% mtransfer_rate+b% mdot_wind_transfer(b% d_i) >= 0 .or. b% CE_flag) then
    b% mdot_system_transfer(b% d_i) = 0d0
    b% mdot_system_transfer(b% a_i) = 0d0
    b% mdot_system_cct = 0d0
else 
    b% mdot_system_transfer(b% d_i) = b% mtransfer_rate * b% mass_transfer_alpha
    b% mdot_system_cct = b% mtransfer_rate * b% mass_transfer_delta
    if (b% point_mass_i == 0 .or. b% model_twins_flag) then
        b% mdot_system_transfer(b% a_i) = b% mtransfer_rate * b% mass_transfer_beta
    else
        ! do not compute mass lost from the vicinity using just mass_transfer_beta, as
        ! mass transfer can be stopped also by going past the Eddington limit
        b% mdot_system_transfer(b% a_i) = (actual_mtransfer_rate + b% component_mdot(b% a_i)) &
            + b% mtransfer_rate * b% mass_transfer_beta
    end if
end if
\end{minted}
This is then used to compute the change in angular momentum in the default routine as:
\begin{minted}[fontsize=\small]{fortran}
! $MESA_DIR/binary/private/binary_jdot.f90
! subroutine default_jdot_ml

!mass lost from circumbinary coplanar toroid
b% jdot_ml = b% jdot_ml + b% mdot_system_cct * b% mass_transfer_gamma * &
     sqrt(standard_cgrav * (b% m(1) + b% m(2)) * b% separation)
\end{minted}

I think just copying this line into my \lstinline[style=output, breaklines=true]|wind_loss| subroutine should be all that is necessary to simulate mass loss from a coplanar circumbinary toroid. The new \lstinline[style=output, breaklines=true]|wind_loss| subroutine is now:
\begin{minted}[fontsize=\small,breaklines]{fortran}
! $RUN_DIR/src/bin/additional_routines.inc

subroutine wind_loss(binary_id, ierr)
    integer, intent(in) :: binary_id
    integer, intent(out) :: ierr
    type(binary_info), pointer :: b
    real(dp) :: q, vw_over_vorb
    real(dp) :: beta, eta
    real(dp) :: omega, I_eff
    real(dp) :: domega_dt_wind

    ierr = 0
    call binary_ptr(binary_id, b, ierr)
    if (ierr /= 0) then
        write (*, *) 'failed in binary_ptr'
        return
    end if

    ! --- compute q == m_accretor / m_donor and v
    q = b%m(b%a_i)/b%m(b%d_i)
    vw_over_vorb = compute_vwind(b%s_donor)/compute_vorbit(b%m(1) + b%m(2), b%separation)
    b%s_donor%xtra(x_vw_over_vorb) = vw_over_vorb

    ! --- get beta and eta from Saladino
    beta = compute_beta_Sal(q, vw_over_vorb)/sqrt(1 - pow2(b%eccentricity))
    b%s_donor%xtra(x_beta) = beta
    eta = compute_eta_Sal(q, vw_over_vorb)
    b%s_donor%xtra(x_eta) = eta

    ! --- star angular momentum loss due to winds
    ! compute the change in spin frequency. this
    ! gets used in extra_binary_finish_step.
    omega = b%s_donor%xtra(x_omega)
    I_eff = b%s_donor%xtra(x_I_eff)
    domega_dt_wind = (b%mdot_system_wind(b%d_i) / (1 - beta) *(b%s_donor%photosphere_r*Rsun)**2*omega/1.5)/I_eff
    b%s_donor%xtra(x_domega_dt_wind) = domega_dt_wind

    ! --- mass lost from vicinity of donor
    b%jdot_ml = b%mdot_system_wind(b%d_i) &
                *eta &
                *(b%separation)**2 &
                *sqrt(1 - pow2(b%eccentricity)) &
                *2*pi/b%period

    ! --- mass lost from circumbinary coplanar toroid
    b% jdot_ml = b% jdot_ml + b% mdot_system_cct * b% mass_transfer_gamma * &
                 sqrt(standard_cgrav * (b% m(1) + b% m(2)) * b% separation)

end subroutine wind_loss
\end{minted}

I simulate an analytic grid again using the same code I used last week to find a nice balance for predicted orbital periods on the MESA grid parameter space. I find that \(\delta =0.2\), \(\gamma =1.25\) gives the following set of predicted orbits:

\begin{figure}[ht]
    \marginfignote{0}{These values give a very broad region in the parameter space that generates periods between 1000 and 5000 days, with a few smaller and larger periods. I think this aligns well with the observed period distribution of barium stars as well, and hopefully should not be \emph{too} challenging for mesa, since we are not trying to create extremely short periods.}
    \centering
    \incplot{w13-analytic-grid-delta}
\end{figure}

\section{Grid}

I ran a grid with the same starting conditions as before, but now using \(\delta =0.2\) and \(\gamma  = 1.25\).

I start by showing the minimum orbit that is achieved during the binary evolution, because this shows a peculiar edge case that I had to account for.
\begin{figure}[ht]
    \marginfignote{0}{Two things should be noted here. \begin{enumd}
    \item  The are significantly more systems that \emph{fail} due to a MESA error, even for large \(q\). 
    \item  There are some systems that reach extremely small periods.
\end{enumd}
}
    \centering
    \incplot{w13-grid-1}
\end{figure}

MESA should prevent systems from evolving past a certain point to stop spending resources on hopeless models. I actually implemented this in the grid already, which is why all models that experience this \emph{issue} have the same minimum period of \(50\) days.

I prevented further evolution by implementing a very simple stopping condition in the \lstinline[style=output, breaklines=true]|extra_binary_finish_step|.

From now on, I quantify these models as unstable as well. The period distribution of the new models looks like this:

\begin{figure}[ht]
    \marginfignote{0}{The periods that we find using with mass loss through the circumbinary ring are \emph{significantly} lower than the periods we find without it. This is obviously what we would expect, but it is nice to see that MESA can, in many cases resolve the RLOF mass transfer, even with this extra mass loss.}
    \centering
    \incplot{w13-grid-2}
\end{figure}

\begin{figure}[ht]
    \marginfignote{0}{This figure compares the orbits from the previous grid with the orbits from the new grid. Clearly, In many cases, the orbits are a lot smaller.}
    \centering
    \incplot{w13-grid-3}
\end{figure}
\nextpage
It is also interesting to look at the maximum ratio of the star radius to the Roche lobe radius.
\begin{figure}[ht]
    \marginfignote{0}{We see two systems with very high ratios, the other systems look more or less like the previous grid. It is easier to again compare the models directly.}
    \centering
    \incplot{w13-grid-4}
\end{figure}


\begin{figure}[ht]
    \marginfignote{0}{Now we see that all models extend past their Roche lobe a bit more when mass is lost from the CCT.}
    \centering
    \incplot{w13-grid-5}
\end{figure}

Now, let's look at the outer Roche lobe radius:
\begin{figure}[ht]
    \marginfignote{0}{All models reach their outer Roche lobe. This is a bit worrying, but it is also how mass is supposed to escape to the CCT.}
    \centering
    \incplot{w13-grid-6}
\end{figure}

Below, I show the models in the fourth column, including the failed model and the very extreme model. This first figure shows the radius of the star and the radius of the Roche lobe as a function of \(q\).
\begin{figure}[ht]
    \centering
    \incplot{w13-extreme1}
\end{figure}

It seems like both the star radius increases with lower initial \(q\), and the Roche lobe radius decreases. That the Roche lobe radius decreases makes sense to me, but I do not know why the star radius increases too.

Another thing to note, is that the orange, green and red model terminate due to low envelope mass, but they are at that point still larger than their Roche lobe. 

In any case, the plot below shows the same plot as above, but now as a fraction.
\begin{figure}[ht]
    \centering
    \incplot{w13-extreme2}
\end{figure}

Now it is very clear that the orange model reaches a much higher maximum than the other models. For comparison, here is the same plot for the fourth column:

\begin{figure}[ht]
    \centering
    \incplot{w13-no-extreme}
\end{figure}

The blue model is one of the models that reached a period smaller than 50 days. It's fraction of \(R_\text{star}/R_\text{RL}\) goes all the way up to about 8 before terminating and has no signs of slowing down.

\nextpage
Lastly, Let's look at the eighth column:

\begin{figure}[ht]
    \centering
    \incplot{w13-extreme4}
\end{figure}

We see a similar effect as in the first set of plots.

\begin{figure}[ht]
    \centering
    \incplot{w13-extreme3}
\end{figure}

Again here. I think the "extreme" points I saw in the 2d grid are look exaggerated due to plotting them on a linear scale and just barely, or just barely not going into the regime where the periods start to drop significantly.

The last thing I want to do is apply the same procedure of finding the final mass ratio to this new set of models.

\begin{figure}[ht]
    \marginfignote{0}{It looks very similar to the plot I showed earlier, meaning the effect of losing some mass to the CCT does not significantly impact the C/O-number ratio.}
    \centering
    \incplot{w13-co-2}
\end{figure}

The plot below shows this more clearly:
\begin{figure}[ht]
    \marginfignote{0}{Here, the two values are divided and what remains looks more or less like noise to me. I think the effects are not super important because the original grid still had some non-conservative mass transfer due to the radiation. This is probably roughly equal in amount to the mass lost to the CCT.}
    \centering
    \incplot{w13-co-3}
\end{figure}

\begin{figure}[ht]
    \marginfignote{0}{This figure shows the differences in mass. Again, they are quite small. Significantly less that 0.8 in any case.}
    \centering
    \incplot{w13-m2}
\end{figure}






\bibliography{master.bib}
\end{document}
